\section*{\centering ВИСНОВКИ}\addcontentsline{toc}{section}{ВИСНОВКИ}

У дисертаційному дослідженні теоретично обґрунтовано, розроблено та експериментально перевірено методику навчання розробки імерсивних освітніх ресурсів майбутніх учителів інформатики. Застосовано змішані методи дослідження (конвергентний паралельний дизайн), що інтегрують кількісний педагогічний експеримент ($N = 64$) та спрямований якісний контент-аналіз відеолекцій (14 транскриптів, 1005 сегментів). Відповідно до поставлених завдань дослідження отримано такі основні результати.

\begin{enumerate}[label=\arabic*.]

\item \textit{Визначено теоретичні засади навчання розробки імерсивних освітніх ресурсів у підготовці майбутніх учителів інформатики} (розділ~1).

Аналіз наукових джерел, структурований за п'ятьма проблемно-тематичними кластерами, засвідчив формування цілісного міждисциплінарного поля досліджень. Метааналітичні дані шести систематичних оглядів підтверджують педагогічну ефективність імерсивних технологій (розміри ефекту $g=0{,}82$--$d=0{,}98$). Встановлено, що когнітивно-афективна модель імерсивного навчання CAMIL з шістьма медіаторними факторами (присутність, агентність, втілення, когнітивне навантаження, мотивація, саморегуляція) та модель TPACK утворюють наскрізну теоретичну рамку дослідження. Аналіз за моделлю TPACK виявив системний дефіцит TPK-компоненту: майбутні вчителі формують TK і PK окремо, проте їх інтеграція в контексті імерсивних середовищ не забезпечується.

Уточнено понятійний апарат: імерсивні освітні ресурси визначено як інтерактивні дидактичні засоби нового покоління; виокремлено п'ять сутнісних характеристик ІОР~-- інтерактивність, імерсивність/присутність, контекстуальність, візуалізація та адаптивність. Обґрунтовано 6-компонентну модель проєктування ІОР (аналіз~$\to$ вимоги~$\to$ вибір ПЗ~$\to$ розробка~$\to$ впровадження~$\to$ оцінювання). Професійну готовність майбутнього вчителя інформатики до розробки ІОР визначено як інтегративну характеристику, що поєднує теоретичні знання, практичні вміння, особистісно-мотиваційні якості та здатність до безперервного професійного зростання.

Порівняльний аналіз ОПП 38~закладів вищої освіти за запропонованою моделлю рівнів зрілості інтеграції ІОР показав, що абсолютна більшість перебуває на рівнях~0--2, тоді як рівні~3--4 представлені одиничними випадками. Дослідницька траєкторія КДПУ (2020--2024) є єдиним задокументованим вітчизняним прикладом системного просування від рівня~2 до рівня~3--4.

\item \textit{Обґрунтовано організаційно-педагогічні основи навчання розробки імерсивних освітніх ресурсів} (розділ~2).

Визначено сукупність одинадцяти дидактичних підходів (аксіологічний, діяльнісний, інноваційний, інтегративний, когнітивний, компетентнісний, конструктивістський, контекстний, рефлексивний, синергетичний, системний), згрупованих у чотири кластери: ціннісно-світоглядний, процесуально-діяльнісний, технологічно-інноваційний та системно-структурний. Побудована матриця відповідності підходів характеристикам ІОР та факторам CAMIL підтверджує повноту запропонованої сукупності. З кластеризованих підходів виведено вісім принципів навчання (науковості, інтерактивності, практичної спрямованості, інтеграції, креативності, варіативності, рефлексивності, безперервності).

Розроблено структурно-функціональну модель підготовки, що охоплює п'ять взаємозумовлених блоків~-- мотиваційно-цільовий, теоретико-методологічний, змістовий, організаційно-технологічний і рефлексивно-оцінний,~-- кожен із яких співвіднесений із конкретними факторами CAMIL та компонентами TPACK. Готовність визначено як інтегративне утворення з шести компонентів (мотиваційно-ціннісний, когнітивний, діяльнісно-проєктувальний, технологічний, інтерактивний, рефлексивно-оцінний).

Визначено чотири педагогічні умови результативності навчання~-- психолого-педагогічну, методичну, технологічну та організаційну,~-- кожна з яких співвіднесена з факторами CAMIL, компонентами TPACK та компонентами готовності. Розроблено триетапну методику навчання розробки ІОР: теоретико-аналітичний етап, практико-розробницький етап (прогресивна послідовність WebAR-проєктів за циклом Колба) та інтегративно-рефлексивний етап. Методику узгоджено з п'ятьма першими принципами навчання Меріла; конкретною її реалізацією виступає авторський курс <<Синтетична реальність в освіті>> (18~модулів) на базі навчального посібника з WebAR.

\item \textit{Організовано та проведено дослідно-експериментальну роботу з перевірки ефективності розробленої методики} (розділ~3).

Розроблено критеріальний апарат: чотири критерії (мотиваційний, когнітивний, операційно-діяльнісний, рефлексивно-оцінювальний) із 24~показниками, чотирма рівнями сформованості та 100-бальною оцінювальною шкалою. Проведено трифазний педагогічний експеримент (2020--2025): констатувальний, формувальний та контрольний етапи за участю 64~студентів ($n_{\text{ЕГ}} = 34$, $n_{\text{КГ}} = 30$).

Результати контрольного етапу засвідчили статистично значущу позитивну динаміку за всіма критеріями: $\chi^2_{\text{емп}} > \chi^2_{\text{крит}} = 7{,}815$ при $\alpha = 0{,}05$. Частка студентів ЕГ з достатнім та високим рівнями зросла від 23,5\% до 70,6\% (КГ: від 23,4\% до 36,7\%). Найбільшу динаміку зафіксовано за операційно-діяльнісним критерієм ($\chi^2_{\text{емп}} = 9{,}217$), що підтверджує ефективність практико-орієнтованого підходу. Апробовано трикрокову методику інтеграції моделей машинного навчання у WebAR (73,5\%~$\to$~58,8\%~$\to$~44,1\% успішності за етапами зростаючої складності).

Спрямований якісний контент-аналіз 14~транскриптів відеолекцій (1005~сегментів, $\approx$23~години) підтвердив реалізацію всіх чотирьох педагогічних умов та виявив темпоральну динаміку факторів CAMIL: мотивація домінує на теоретико-аналітичному етапі, тілесність~-- на практико-розвивальному (пік при вивченні Face Tracking), саморегуляція концентрується на перехідних тижнях. Аналіз стратегій викладання виявив домінування моделювання ($\approx$43\%) із систематичним переходом від інструктор-центрованих до студент-центрованих стратегій. Тріангуляція кількісних та якісних даних (конвергентний паралельний дизайн) продемонструвала збіжність за всіма критеріями; часткова конвергенція за рефлексивно-оцінювальним критерієм пояснюється тим, що саморегуляція формується переважно через письмові форми.

\end{enumerate}

Таким чином, гіпотезу дослідження підтверджено: методика навчання розробки імерсивних освітніх ресурсів є ефективною за умови реалізації обґрунтованих педагогічних умов, що підтверджено як статистично значущими кількісними результатами, так і якісними свідченнями педагогічних механізмів. Інтеграція кількісних та якісних методів аналізу дала змогу не лише зафіксувати \textit{що} змінилося, а й розкрити \textit{як} саме навчальний процес забезпечив ці зміни.

Перспективними напрямами подальших досліджень є: розширення методики на підготовку вчителів інших спеціальностей; інтеграція генеративного штучного інтелекту в процес розробки ІОР; лонгітюдне дослідження впливу навчання розробки ІОР на педагогічну практику випускників; масштабування методики для дистанційного та змішаного навчання; поглиблення якісного аналізу студентських артефактів та рефлексій як додаткове джерело тріангуляції.

\label{end_main_text}
